% GENERATED FILE. Edit canonical lessons, then run npm run generate:books.
% canonical-source-hash: fnv1a64:8d76ee08b6b2b43f
% canonical-lessons: TE-C202-carcincu, TE-C202-sampradincu, TE-C202-tanikhi-ceyu, TE-C202-namodu-ceyu, TE-C202-sadhana-ceyu

\chapter{To discuss, To contact, To check, To register, To practise}
\label{ch:te-to-discuss-to-contact-to-check-to-register-to-practise}
\hlchaptermodality{202}

\begin{chapteropening}
\textbf{By the end of this chapter:} \emph{I can say to discuss, to contact, to check, to register and to practise.}

Complete the last lesson of chapter 202: I can say to discuss, to contact, to check, to register and to practise.
\end{chapteropening}

\section[\texorpdfstring{\te{చర్చించు} (carciñcu)}{చర్చించు (carciñcu)}]{\te{చర్చించు} (carciñcu) — to discuss}

\label{lesson:TE-C202-carcincu}



\begin{quote}
Before the new one: say the Telugu for to shiver, then the Telugu for to move.
\end{quote}

\subsection*{You'll want to know: \te{చర్చించు}}
\textbf{\te{చర్చించు}} — \emph{carciñcu} — \textquotedblleft{}to discuss\textquotedblright{}.

\begin{grammarlens}[title={What you've built}]
One more word for this chapter.
\end{grammarlens}

\subsection*{Your turn}
\begin{itemize}
\raggedright
  \item \emph{Say it:} \emph{carciñcu}
  \item \emph{Say it:} \emph{carciñcu}, once more
  \item \emph{Say it:} say \emph{carciñcu}
  \item \emph{Recall:} say the Telugu for to shiver, then the Telugu for to move, then say \emph{carciñcu} again
\end{itemize}

\begin{tcolorbox}[breakable,colback=teal!4,colframe=teal!35!black,title={Before you move on}]
What does \te{చర్చించు} mean? (\textquotedblleft{}To discuss\textquotedblright{}.) Say it once more.
\end{tcolorbox}
\section[\texorpdfstring{\te{సంప్రదించు} (sampradiñcu)}{సంప్రదించు (sampradiñcu)}]{\te{సంప్రదించు} (sampradiñcu) — to contact, to consult}

\label{lesson:TE-C202-sampradincu}



\begin{quote}
Before the new one: say the Telugu for to move, then the Telugu for to discuss.
\end{quote}

\subsection*{You'll want to know: \te{సంప్రదించు}}
\textbf{\te{సంప్రదించు}} — \emph{sampradiñcu} — \textquotedblleft{}to contact, to consult\textquotedblright{}.

\begin{grammarlens}[title={What you've built}]
One more word for this chapter.
\end{grammarlens}

\subsection*{Your turn}
\begin{itemize}
\raggedright
  \item \emph{Say it:} \emph{sampradiñcu}
  \item \emph{Say it:} \emph{sampradiñcu}, once more
  \item \emph{Say it:} say \emph{sampradiñcu}
  \item \emph{Recall:} say the Telugu for to move, then the Telugu for to discuss, then say \emph{sampradiñcu} again
\end{itemize}

\begin{tcolorbox}[breakable,colback=teal!4,colframe=teal!35!black,title={Before you move on}]
What does \te{సంప్రదించు} mean? (\textquotedblleft{}To contact, to consult\textquotedblright{}.) Say it once more.
\end{tcolorbox}
\section[\texorpdfstring{\te{తనిఖీ} \te{చేయు} (tanikhī cēyu)}{తనిఖీ చేయు (tanikhī cēyu)}]{\te{తనిఖీ} \te{చేయు} (tanikhī cēyu) — to check, to inspect}

\label{lesson:TE-C202-tanikhi-ceyu}



\begin{quote}
Before the new one: say the Telugu for to discuss, then the Telugu for to contact.
\end{quote}

\subsection*{You'll want to know: \te{తనిఖీ} \te{చేయు}}
\textbf{\te{తనిఖీ} \te{చేయు}} — \emph{tanikhī cēyu} — \textquotedblleft{}to check, to inspect\textquotedblright{}.

\begin{grammarlens}[title={What you've built}]
One more word for this chapter.
\end{grammarlens}

\subsection*{Your turn}
\begin{itemize}
\raggedright
  \item \emph{Say it:} \emph{tanikhī cēyu}
  \item \emph{Say it:} \emph{tanikhī cēyu}, once more
  \item \emph{Say it:} say \emph{tanikhī cēyu}
  \item \emph{Recall:} say the Telugu for to discuss, then the Telugu for to contact, then say \emph{tanikhī cēyu} again
\end{itemize}

\begin{tcolorbox}[breakable,colback=teal!4,colframe=teal!35!black,title={Before you move on}]
What does \te{తనిఖీ} \te{చేయు} mean? (\textquotedblleft{}To check, to inspect\textquotedblright{}.) Say it once more.
\end{tcolorbox}
\section[\texorpdfstring{\te{నమోదు} \te{చేయు} (namōdu cēyu)}{నమోదు చేయు (namōdu cēyu)}]{\te{నమోదు} \te{చేయు} (namōdu cēyu) — to register, to record}

\label{lesson:TE-C202-namodu-ceyu}



\begin{quote}
Before the new one: say the Telugu for to contact, then the Telugu for to check.
\end{quote}

\subsection*{You'll want to know: \te{నమోదు} \te{చేయు}}
\textbf{\te{నమోదు} \te{చేయు}} — \emph{namōdu cēyu} — \textquotedblleft{}to register, to record\textquotedblright{}.

\begin{grammarlens}[title={What you've built}]
One more word for this chapter.
\end{grammarlens}

\subsection*{Your turn}
\begin{itemize}
\raggedright
  \item \emph{Say it:} \emph{namōdu cēyu}
  \item \emph{Say it:} \emph{namōdu cēyu}, once more
  \item \emph{Say it:} say \emph{namōdu cēyu}
  \item \emph{Recall:} say the Telugu for to contact, then the Telugu for to check, then say \emph{namōdu cēyu} again
\end{itemize}

\begin{tcolorbox}[breakable,colback=teal!4,colframe=teal!35!black,title={Before you move on}]
What does \te{నమోదు} \te{చేయు} mean? (\textquotedblleft{}To register, to record\textquotedblright{}.) Say it once more.
\end{tcolorbox}
\section[\texorpdfstring{\te{సాధన} \te{చేయు} (sādhana cēyu)}{సాధన చేయు (sādhana cēyu)}]{\te{సాధన} \te{చేయు} (sādhana cēyu) — to practise}

\label{lesson:TE-C202-sadhana-ceyu}



\begin{quote}
Before the new one: say the Telugu for to check, then the Telugu for to register.
\end{quote}

\subsection*{You'll want to know: \te{సాధన} \te{చేయు}}
\textbf{\te{సాధన} \te{చేయు}} — \emph{sādhana cēyu} — \textquotedblleft{}to practise\textquotedblright{}.

\begin{grammarlens}[title={What you've built}]
One more word for this chapter.
\end{grammarlens}

\subsection*{Your turn}
\begin{itemize}
\raggedright
  \item \emph{Say it:} \emph{sādhana cēyu}
  \item \emph{Say it:} \emph{sādhana cēyu}, once more
  \item \emph{Say it:} say \emph{sādhana cēyu}
  \item \emph{Recall:} say the Telugu for to check, then the Telugu for to register, then say \emph{sādhana cēyu} again
\end{itemize}

\begin{tcolorbox}[breakable,colback=teal!4,colframe=teal!35!black,title={Before you move on}]
What does \te{సాధన} \te{చేయు} mean? (\textquotedblleft{}To practise\textquotedblright{}.) Say it once more.
\end{tcolorbox}
